Los experimentos se ejecutaron en un equipo con procesador AMD Ryzen 7
5700X, 24\,GB de RAM a 2133\,MT/s, bajo Windows 11 (versión 25H2, compilación
26200.9445). El código se compiló con \texttt{g++} 16.2.0 (MSYS2,
Rev3) usando el estándar C++17 y el flag de optimización \texttt{-O2}
(\texttt{g++ -O2 -std=c++17}), consistente con el \texttt{makefile}
incluido en el repositorio.

Los datasets utilizados corresponden a los descritos en el enunciado de
la tarea: para ordenamiento, arreglos de $n \in \{10, 10^3, 10^5, 10^7\}$
elementos, con tipo $\in \{\text{ascendente}, \text{descendente},
\text{aleatorio}\}$ y dominio $\in \{D1, D7\}$; para multiplicación de
matrices, matrices cuadradas de $n \in \{16, 64, 256\}$\footnote{Se
excluyó $n=1024$ ($2^{10}$) de las mediciones, con autorización del
curso.}, con tipo $\in \{\text{densa}, \text{diagonal}, \text{dispersa}\}$
y dominio $\in \{D0, D10\}$. En ambos casos se generaron 3 muestras
aleatorias independientes por combinación de parámetros.